\section[Conclusion]{Conclusion}

\begin{frame}{Bilan technique des optimisations}
    \begin{block}{Synthèse des gains en C++}
        \fontsize{8}{9}\selectfont
        \begin{itemize}
            \item \textbf{Levée des verrous} : L'identification de l'aliasing comme goulot d'étranglement a permis une restructuration sémantique efficace (\texttt{\_\_restrict}).
            \item \textbf{Efficacité des noyaux} : La charge relative de la routine \texttt{resid} a été réduite de \textbf{47,85 \%} à \textbf{29,57 \%} sur AMD.
            \item \textbf{Paradoxe MAQAO} : Le score de compilation reste à \textbf{16,7 \%}, confirmant un défaut de détection des métadonnées du binaire C++, ce qui était gênant pour nos déductions initiales.
        \end{itemize}
    \end{block}
\end{frame}



\begin{frame}{Le Duel Final : C++ Optimisé vs Fortran Natif}
    \centering
    \inputImage{figures/comparison_duel_report_amd_cpp_vs_f90_opt.png}{Comparaison finale sur AMD : Fortran vs C++ optimisé}{0.8\textwidth}{!}{fig:comparison_duel_report_amd_cpp_vs_f90_opt}
    
    \vspace{-0.5cm}
    \begin{columns}[T]
        \begin{column}{0.48\textwidth}
            \begin{exampleblock}{\footnotesize Fortran (Référence)}
                \fontsize{7}{8}\selectfont
                \begin{itemize}
                    \item Temps total : \textbf{11,83 s}.
                    \item Score : \textbf{100 \%}.
                    \item Ratio d'activité : \textbf{99,6 \%}.
                \end{itemize}
            \end{exampleblock}
        \end{column}
        \begin{column}{0.48\textwidth}
            \begin{alertblock}{\footnotesize C++ (Optimisé)}
                \fontsize{7}{8}\selectfont
                \begin{itemize}
                    \item Temps total : \textbf{19,48 s}.
                    \item Score : \textbf{16,7 \%}.
                    \item Ratio d'activité : \textbf{64,1 \%}.
                \end{itemize}
            \end{alertblock}
        \end{column}
    \end{columns}
\end{frame}

\begin{frame}{Conclusion}
    \begin{itemize}
        \fontsize{9}{10}\selectfont
        \item \textbf{Performance} : Malgré les optimisations manuelles, le Fortran reste \textbf{1,65x plus rapide} que le C++ sur cet algorithme multigrid.
        \item \textbf{Architecture} : Le Fortran exploite nativement mieux le parallélisme (8 threads actifs contre 5 pour le C++), avec un ratio d'activité quasi parfait (99,6 \%).
    \end{itemize}
    
    \begin{block}{}
        \centering \fontsize{12}{14}\selectfont
        \textit{Le Fortran demeure virtuellement imbattable pour les codes de physique numérique. Sa gestion native de l'absence d'aliasing et sa structure orientée tableaux lui permettent d'atteindre 100 \% de score d'optimisation là où le C++ se heurte à des contraintes sémantiques complexes.
        On constate ainsi la lourdeur d'optimiser plus loin le C++ pour tenter d'égaler le Fortran sur certains algorithmes de problèmes physiques.}
    \end{block}
\end{frame}



\begin{frame}{Perspectives}
    \vspace{0.6cm}
    \centering
    \textbf{\huge Merci de votre attention} \\
    \vspace{0.3cm}
    \textit{Avez-vous des questions ?}
\end{frame}